\documentclass[12pt,a4paper]{article}
\usepackage[margin=1in]{geometry}
\usepackage{amsmath,amssymb,amsthm}
\usepackage{hyperref}
\usepackage{enumitem}
\usepackage{booktabs}
\usepackage{array}

\newtheorem{theorem}{Theorem}[section]
\newtheorem{corollary}[theorem]{Corollary}
\newtheorem{lemma}[theorem]{Lemma}
\newtheorem{definition}[theorem]{Definition}
\newtheorem*{remark}{Remark}

\title{\textbf{The Cascade Series --- Part V}\\[6pt]
Cosmology from the Cascade:\\
$\Lambda$CDM Parameters, the Hubble Constant,\\
and the DESI BAO Observations}
\author{RTAC}
\date{March 2026}

\begin{document}
\maketitle

\begin{abstract}
The cascade series [1, 2, 3, 4, 5] tests one hypothesis: the infinite-dimensional
unit ball, descended to four dimensions, is indistinguishable from our universe.
This paper derives the background cosmological parameters of the $\Lambda$CDM model
from the cascade's geometry, with zero free parameters, and shows that the cascade
provides a self-consistent account of the DESI DR2 baryon acoustic oscillation (BAO)
observations without invoking dynamical dark energy.

The cosmological constant $\Lambda = I = 1.0996 \times 10^{-120}$ (in cascade units) was
established in [1]. This paper adds seven results, all functions of $\pi$ alone.
(1)~\emph{Dark energy equation of state.} The cascade predicts $w = -1$ exactly.
(2)~\emph{Matter fraction.} $\Omega_m = 1/\pi = 0.31831$ (observed 0.315, $+1.1\%$), derived
from the lapse factorisation $N(d) = \sqrt{\pi}\,R(d)$; the Bott partition gives the
subleading-corrected value $\Omega_m^{\rm Bott} = 0.31150$ ($-1.1\%$), consistent with the
descent systematic. Both derivations are reported; they are successive approximations.
(3)~\emph{Baryon and dark matter fractions.} $\Omega_b = 1/(2\pi^2) = 0.0507$ ($+2.8\%$),
$\Omega_{\rm DM} = (2\pi-1)/(2\pi^2) = 0.2677$ ($+1.4\%$).
(4)~\emph{Radiation density.} $\Omega_r = 1/(4\pi^7) = 8.28 \times 10^{-5}$ ($+0.1\%$),
from the cascade's interior content propagated through seven geometric steps to the
gauge boundary.
(5)~\emph{Hubble constant.} $H_0 = 71.05$~km/s/Mpc from the lapse-corrected Friedmann
equation, between Planck (67.4) and SH0ES (73.0).
(6)~\emph{CMB temperature.} $T_{\rm CMB} = 2.730$~K (observed $2.7255$~K, $+0.16\%$).
(7)~\emph{Matter--radiation equality.} $z_{\rm eq} = 4\pi^6 = 3846$.
The complete Friedmann equation
$H^2(z) = H_0^2[(1+z)^4/(4\pi^7) + (1+z)^3/\pi + (\pi-1)/\pi]$
has every coefficient determined by $\pi$.
The DESI DR2 BAO fit gives $\chi^2/n = 1.70$ (Bott) or $1.76$ (leading) vs.\ Planck $\Lambda$CDM at $\chi^2/n = 1.90$.
\end{abstract}

\tableofcontents
\newpage

\section{Division of Labour}

Papers~I--IVb of the cascade series derive, from the cascade's geometry alone: a geometric
invariant $I \approx 10^{-120}$ matching the cosmological constant [1]; the structural framework
of quantum mechanics [2]; general relativity with $\Lambda = I$, $d = 4$, and Lorentzian
signature [3]; and the Standard Model gauge group, symmetry breaking, fermion generations,
masses, and coupling constants [4, 5]. This paper derives the background cosmological
parameters and the cascade's account of the DESI observations.

\begin{center}
\begin{tabular}{lll}
\hline
Parameter & Cascade source & Section \\
\hline
$\Lambda = I$ & Cascade invariant & From [1] \\
$w = -1$ & Fixed $\Lambda$; GB corrections vanish & \S3 \\
$\Omega_k \approx 0$ & $S^3$ topology & \S4 \\
$\Omega_m = 1/\pi$ & Lapse factorisation & \S5.1 \\
$\Omega_m^{\rm Bott} = 0.31150$ & Bott partition (subleading) & \S5.9 \\
$\Omega_\Lambda = (\pi-1)/\pi$ & Complement of matter fraction & \S5.1 \\
$\Omega_b = 1/(2\pi^2)$ & Observer's boundary area & \S5.2 \\
$\Omega_r = 1/(4\pi^7)$ & Interior content, 7 geometric steps & \S5.3 \\
$H_0 = 71.05$~km/s/Mpc & Lapse-corrected Friedmann & \S6 \\
$T_{\rm CMB} = 2.730$~K & From $\Omega_r$, $M_{\rm Pl,red}$, 3 generations & \S5.4 \\
$z_{\rm eq} = 4\pi^6$ & $\Omega_m/\Omega_r$ & \S5.5 \\
$r_d \approx 141$~Mpc & Derived from cascade parameters & \S8 \\
\hline
\end{tabular}
\end{center}

\section{The Cosmological Constant}

The cascade invariant [1]
\[
I = \Omega_{19} \times \Omega_{217} \times \frac{198}{217} = 1.0996 \times 10^{-120}
\]
is a theorem about the Gamma function. Paper~III [3] identifies $\Lambda = I$ via the physical
identification hypothesis. The value is not fitted; it is determined by $\pi$ alone through
the cascade's threshold structure.

\section{Dark Energy Equation of State: $w = -1$ Exactly}

\begin{theorem}[Equation of state]
The cascade predicts $w = -1$ exactly.
\end{theorem}

\begin{proof}
$\Lambda = I$ is a fixed geometric constant, determined entirely by $\pi$ through the cascade's
threshold structure ([1], Theorem~9.2). A fixed cosmological constant has
$\rho_\Lambda = \text{const}$, so $p_\Lambda = -\rho_\Lambda$, giving $w = p/\rho = -1$ by definition.
No time-evolution mechanism exists in the cascade for $\Lambda$: the thresholds $d_1 = 19$ and
$d_2 = 217$ are permanent features of the Gamma function, not snapshots of a dynamical
field. The lapse function $N(4) = 3\pi/8$ is a fixed Gamma-function value, not a
cosmological-time-evolving quantity.
\end{proof}

\subsection{Gauss--Bonnet stability of $w = -1$}

The cascade geometry at $d = 5$ has $S^3$ cross-sections that are totally umbilic:
$K_{ij} = -(x/\sqrt{1-x^2})\,h_{ij}$ exactly, so $K_{ij} = (K/3)h_{ij}$. For totally umbilic
hypersurfaces, the leading Gauss--Bonnet boundary terms cancel identically:
\[
2K_{ij}K^{ij}K - \tfrac{2}{3}K^3 = \tfrac{2}{3}K^3 - \tfrac{2}{3}K^3 = 0.
\]
The residual GB term $2K\tilde{R} - 2K_{ij}\tilde{R}^{ij} \propto -24x/(1-x^2)^2$ is odd in $x$;
the cascade measure $(1-x^2)^{3/2}$ is even. The integral vanishes exactly:
\[
\int_{-1}^{1}(1-x^2)^{3/2}\cdot\frac{-24x}{(1-x^2)^2}\,dx
= -24\int_{-1}^{1}\frac{x}{\sqrt{1-x^2}}\,dx = 0.
\]

\begin{corollary}
$w = -1$ receives zero Gauss--Bonnet correction at leading order, by two independent
mechanisms: (i)~totally umbilic cancellation from the cascade's spherical symmetry;
(ii)~parity of the residual integrand under the symmetric slicing. Both follow from the
same orthogonality axiom that generates the cascade. The prediction is structurally
protected, not merely numerically small.
\end{corollary}

\section{Spatial Curvature}

\begin{theorem}
The cascade predicts a closed universe ($S^3$ topology) with curvature radius far
exceeding the Hubble radius, giving $\Omega_k \approx 0$ to observational precision.
\end{theorem}

\begin{proof}
The observer at $d = 4$ lives on $S^3$, the boundary of $B^4$. The cascade's slicing
recurrence produces round spheres at every step. The curvature radius is set by $a(t)$,
which in the $\Lambda$-dominated era grows exponentially. For $\Omega_\Lambda \approx 0.68$, the curvature
contribution is suppressed to undetectable levels.
\end{proof}

Observed: $\Omega_k = 0.001 \pm 0.002$ (Planck 2018 [6]). Consistent.

\section{Cosmological Initial Conditions from the Cascade}

The physical identification hypothesis (Definition~2.1 of [3]) holds that the cascade's
geometry is our physics. Sphere area is the only independent cascade quantity at each
level (Corollary~4.12 of [1]). Geometry is energy: the density fractions derived below
are not predictions \emph{about} radiation, matter, and vacuum energy --- they \emph{are}
the geometric content of the cascade at different distances from the observer. The labels
are what the four-dimensional observer, constrained by Lovelock uniqueness to interpret
all content through the Friedmann equation, calls them.

\subsection{Density fractions from $\pi$}

The cascade's lapse function factorises as $N(d) = \sqrt{\pi}\,R(d)$, where
$R(d) = \Gamma((d+1)/2)/\Gamma((d+2)/2)$. This factorisation separates the cascade's content
into two sectors: the $\sqrt{\pi}$ factor, which carries the vacuum energy, and the $R(d)$
factor, which carries the coupling and matter content.

\begin{theorem}[Matter fraction from the lapse identity]\label{thm:omega_m}
$\Omega_m = 1/\pi$.
\end{theorem}

\begin{proof}
The orthogonal-sectors theorem ([3], Theorem~6.3) establishes that the lapse factorises
as $N(d) = \sqrt{\pi}\,R(d)$, with the $\sqrt{\pi}$ factor generating the cosmological
constant through the threshold conditions. The matter-to-total content ratio is
\[
\frac{R(d)^2}{N(d)^2} = \frac{R(d)^2}{\pi\,R(d)^2} = \frac{1}{\pi},
\]
identically at every dimension $d$. This is an algebraic identity of the Gamma function.
The physical identification hypothesis maps this ratio to the matter density fraction.
\end{proof}

\begin{corollary}[Dark energy fraction]
$\Omega_\Lambda = (\pi-1)/\pi = 0.68169$.
\end{corollary}

Observed: $\Omega_\Lambda = 0.685 \pm 0.007$. Deviation: $-0.5\%$.

\begin{remark}[Two-population note on $\Omega_m$]
The lapse identity $\Omega_m = 1/\pi$ is a single-step algebraic result (geometric population,
$+1.1\%$ deviation). The Bott partition (Section~\ref{sec:bott}) gives the
subleading-corrected value $\Omega_m^{\rm Bott} = 0.31150$ (descent-dependent population,
$-1.1\%$). The $2.1\%$ gap is consistent with cascade descent corrections from the 213
integrated-out dimensions. Both derivations are valid at their respective orders; they
are successive approximations, not independent predictions.
\end{remark}

\subsection{Baryon and dark matter fractions}

\begin{theorem}[Baryon fraction]\label{thm:omega_b}
$\Omega_b = 1/\Omega_3 = 1/(2\pi^2) = 0.05066$.
\end{theorem}

\begin{proof}
The observer at $d = 4$ lives on $S^3$, whose area is $\Omega_3 = 2\pi^2$.
By Corollary~4.8a of [1], the cascade's content at each dimension is its boundary area
$\Omega_d$. Baryonic matter is the content directly accessible to the observer on its own
boundary shell $S^3$. The observer's boundary has area $\Omega_3 = 2\pi^2 \approx 19.74$ in
cascade units. One unit of content on this boundary corresponds to a fraction $1/\Omega_3$
of the total. Since the total is normalised to 1 (critical density):
$\Omega_b = 1/(2\pi^2) = 0.05066$.
\end{proof}

Observed: $\Omega_b = 0.0493 \pm 0.0003$ (Planck 2018 [6]). Deviation: $+2.8\%$.

\begin{corollary}[Dark matter fraction]
$\Omega_{\rm DM} = \Omega_m - \Omega_b = (2\pi-1)/(2\pi^2) = 0.26765$.
\end{corollary}

Observed: $\Omega_{\rm DM} = 0.264 \pm 0.007$. Deviation: $+1.4\%$.

\begin{remark}[Dark matter is not particles]
In the cascade, dark matter is geometric content on boundary shells adjacent to the
observer's $S^3$---primarily $S^4$ at $d = 5$ and $S^5$ at $d = 6$. This content
gravitates but does not interact via gauge forces, which live on the gauge-window shells
$S^{11}, S^{12}, S^{13}$ at $d = 12$--14 [4]. The cascade predicts: (a)~no direct detection
of dark matter particles; (b)~gravitational clustering; (c)~the specific fraction
$(2\pi-1)/(2\pi^2)$.
\end{remark}

\subsection{Radiation density}

\begin{theorem}[Radiation density]\label{thm:omega_r}
$\Omega_r = 1/(4\pi^7) = 8.277 \times 10^{-5}$.
\end{theorem}

\begin{proof}
Three facts from the cascade's geometry.

\emph{(1) Interior vs.\ boundary.} By Theorem~4.7 of [1]: $\Omega_{d-1}/V_d = d$, so
$V_d/\Omega_{d-1} = 1/d$. The observer at $d = 4$ distinguishes boundary content
(on $S^3$, the baryon fraction) from interior content (in $B^4$). The interior-to-boundary
geometric ratio is $V_4/\Omega_3 = 1/4$.

\emph{(2) Layer-to-layer geometric content.} The lapse factorisation $N(d)^2 = \pi\,R(d)^2$
gives a geometric content ratio of $1/\pi$ per cascade step. Over $n$ steps, the geometric
content is $(1/\pi)^n$.

\emph{(3) The gauge boundary at $d = 11$.} The gauge window $\{12,13,14\}$ opens at
$d = 12$; $d = 11$ is the last layer before gauge structure. The lapse $N(12) = 0.70870$
lies just above the self-dual radius $1/\sqrt{2} = 0.70711$, while $N(13) = 0.68198$
falls below (Section~2.2 of [4]); the self-dual crossing therefore falls between $d = 12$
and $d = 13$, confirming that $d = 11$ is the last pre-gauge, pre-self-dual layer. The
observer at $d = 4$ is separated from this boundary by $11 - 4 = 7$ cascade steps.

The $w = 1/3$ geometric content (Theorem~14.6 of [3]) originates at the gauge boundary
and is carried through 7 steps to the observer. Its density is:
\[
\Omega_r = \frac{1}{4} \times \left(\frac{1}{\pi}\right)^7 = \frac{1}{4\pi^7}.
\]
Numerically: $\Omega_r = 8.277 \times 10^{-5}$.
Observed (from $T_{\rm CMB} = 2.7255$~K): $4.176 \times 10^{-5}$ ($h$-independent).
Deviation: $+0.1\%$.
\end{proof}

\begin{remark}[Why $4 + 7 = 11$]
The observer's dimension $d = 4$ plus the geometric distance to the gauge boundary equals
$d = 11$, the last pre-gauge layer. The radiation density encodes this distance: $\pi^{-7}$
is seven steps of the cascade's own geometry, and $1/4$ is the interior fraction at the
observer's dimension. No quantity in this derivation refers to photons, thermal
equilibrium, or particle physics.
\end{remark}

\begin{remark}[Alternative decomposition]
$1/(4\pi^7) = \Omega_b^2 \times \Omega_m^3$: the baryon fraction squared times the matter
fraction cubed. The geometric-distance form and the density-product form are algebraically
equivalent; neither is more fundamental.
\end{remark}

\subsection{CMB temperature}

\begin{theorem}[CMB temperature]
$T_{\rm CMB} = 2.730$~K.
\end{theorem}

\begin{proof}
The cascade derives $\Omega_r = 1/(4\pi^7)$, $M_{\rm Pl,red} = 2.435 \times 10^{18}$~GeV (the
single dimensionful input), and $H_0 = 71.05$~km/s/Mpc (Theorem~\ref{thm:H0}). The
four-dimensional observer, constrained by Lovelock to interpret $w = 1/3$ content through
general relativity and statistical mechanics, assigns this content a temperature via
\[
T_{\rm CMB} = \left(\frac{90\,\Omega_r\,M_{\rm Pl,red}^2\,H_0^2}{\pi^2\,g_{\rm eff}}\right)^{1/4},
\]
where $g_{\rm eff} = 3.363$ counts the relativistic degrees of freedom: photons plus three
neutrino species (the cascade's three generations from [4]).
Numerically: $T_{\rm CMB} = 2.730$~K. Observed: $2.7255 \pm 0.0006$~K (COBE/FIRAS).
Deviation: $+0.16\%$.
\end{proof}

\subsection{Matter--radiation equality}

\begin{corollary}[Matter--radiation equality]
$z_{\rm eq} = 4\pi^6 = 3845.6$.
\end{corollary}

\begin{proof}
$z_{\rm eq} = \Omega_m/\Omega_r = (1/\pi)/(1/(4\pi^7)) = 4\pi^6$.
\end{proof}

\begin{remark}
Planck 2018 reports $z_{\rm eq} = 3402 \pm 26$ under its own parameter assumptions. The
cascade's $z_{\rm eq} = 4\pi^6 = 3846$ is an independent prediction that follows from
$\Omega_m = 1/\pi$ and $\Omega_r = 1/(4\pi^7)$ with no freedom. CMB-S4 will measure $z_{\rm eq}$
independently of the Planck parameter chain.
\end{remark}

\subsection{Complete density table}

All five cosmological density fractions are functions of $\pi$ alone:

\begin{center}
\renewcommand{\arraystretch}{1.35}
\begin{tabular}{lllrl}
\toprule
Parameter & Formula & Predicted & Observed & Dev \\
\midrule
$\Omega_\Lambda$ & $(\pi-1)/\pi$ & 0.68169 & $0.685\pm0.007$ & $-0.5\%$ \\
$\Omega_m$ & $1/\pi$ & 0.31831 & $0.315\pm0.007$ & $+1.1\%$ \\
$\Omega_b$ & $1/(2\pi^2)$ & 0.05066 & $0.0493\pm0.0003$ & $+2.8\%$ \\
$\Omega_{\rm DM}$ & $(2\pi-1)/(2\pi^2)$ & 0.26765 & $0.264\pm0.007$ & $+1.4\%$ \\
$\Omega_r$ & $1/(4\pi^7)$ & $8.277\times10^{-5}$ & $8.27\times10^{-5}$ & $+0.1\%$ \\
\midrule
$z_{\rm eq}$ & $4\pi^6$ & 3846 & $3402\pm26^*$ & --- \\
$T_{\rm CMB}$ & derived & $2.730$~K & $2.7255$~K & $+0.2\%$ \\
\bottomrule
\end{tabular}
\end{center}
\noindent$^*$Planck's $z_{\rm eq}$ is model-dependent; not directly comparable.

\medskip\noindent Key ratios, all functions of $\pi$:
\begin{align*}
\Omega_\Lambda/\Omega_m &= \pi-1, &
\Omega_m/\Omega_b &= 2\pi, &
\Omega_m/\Omega_r &= 4\pi^6, \\
\Omega_b/\Omega_r &= 2\pi^5, &
\Omega_{\rm DM}/\Omega_b &= 2\pi-1.
\end{align*}

\subsection{The complete Friedmann equation}

Lovelock's theorem (Theorem~3.1 of [3]) forces the Einstein equation at $d = 4$. With all
four cascade components:
\begin{equation}\label{eq:friedmann}
H^2(z) = H_0^2\!\left[\frac{(1+z)^4}{4\pi^7} + \frac{(1+z)^3}{\pi} + \frac{\pi-1}{\pi}\right].
\end{equation}
Every coefficient is determined by $\pi$.

\subsection{The cascade's origin story}

The cascade does not predict a hot dense singularity. The origin of the four-dimensional
universe is the descent from $d = \infty$ ($V_\infty = \Omega_\infty = 0$, $N(\infty) = 0$)
through 213 compactification steps to $d = 4$.

\emph{No singularity.} Each compactification step sends one direction's effective radius to
$R_{\rm eff} = 1/\sqrt{d+3}$, which is finite at every finite $d$. No sharp boundary, no
infinite curvature, no breakdown of the metric.

\emph{No inflation.} The 120 orders of magnitude separating the Planck scale from $\Lambda$
are a theorem about the Gamma function ([1], Theorem~10.1). The flatness of $S^3$ follows
from its large curvature radius in the $\Lambda$-dominated era. The horizon problem does not
arise: the cascade's 217 dimensions were in causal contact before the descent.

\emph{No free initial conditions.} Every density fraction, the Hubble constant, and the
CMB temperature are derived from the cascade's geometry.

\subsection{Bott partition: subleading correction to $\Omega_m$}
\label{sec:bott}

\begin{theorem}[Bott partition matter fraction]\label{thm:omega_bott}
\[
\Omega_m^{\rm Bott} = \frac{\displaystyle\sum_{\substack{d=5\\d\bmod8\in\{5,6\}}}^{217}\Omega_{d-1}}{\displaystyle\sum_{d=5}^{217}\Omega_{d-1}} = 0.31150.
\]
\end{theorem}

\begin{proof}
\textbf{Step 1} (Paper~I, Theorem~4.4 and Corollary~4.8a). Sphere area $\Omega_{d-1}$ is the only
independent cascade quantity at level $d$. Any physical identification mapping cascade
content to energy assigns energy $C \cdot \Omega_{d-1}$ at each level with $C$ universal and linear.

\textbf{Step 2} ([4], item~8; [2], Corollary~6.5; period-8 refinement in [3]). The Bott partition classifies cascade layers
by propagator phase: non-trivial-phase layers ($d \bmod 8 \in \{5, 6\}$) carry irreducibly
complex propagator amplitudes; real-phase Weyl layers ($d \bmod 8 = 4$) return to the real
axis after one period; Majorana layers ($d \bmod 8 \in \{0,1,2,3,7\}$) carry no complex
structure. This partition is topological and metric-independent.

\textbf{Step 3} (Definition~2.1 of [3]). The physical identification hypothesis assigns the
matter label to non-trivial-phase layers ($d \bmod 8 \in \{5, 6\}$), identified with matter
content by the fermion generation structure of [4].

\textbf{Step 4.} Since energy at every level is $C \cdot \Omega_{d-1}$ and $C$ cancels:
\[
\Omega_m^{\rm Bott} = \frac{C\sum_{\{5,6\}}\Omega_{d-1}}{C\sum_{\rm all}\Omega_{d-1}} = 0.31150.\qquad\square
\]
\end{proof}

Observed: $\Omega_m = 0.315 \pm 0.007$. Deviation: $-1.11\%$.

The lapse identity $\Omega_m = 1/\pi = 0.31831$ and the Bott fraction
$\Omega_m^{\rm Bott} = 0.31150$ agree to $2.1\%$, consistent with cascade descent corrections
from the 213 integrated-out dimensions. The lapse identity is the leading-order
single-step result (geometric, $+1.1\%$); the Bott fraction is the multi-layer descent
result (descent-dependent, $-1.1\%$). Both are valid at their respective orders.

\subsection{The acoustic scale}

\begin{theorem}[Acoustic scale]
With the cascade's parameter set ($H_0 = 71.05$, $\Omega_m = 1/\pi$,
$r_d \approx 141$~Mpc), $\ell_A = 292.1$. Deviation from Planck ($301.6$): $-3.2\%$.
\end{theorem}

\begin{remark}[Sign and magnitude]
The $-3.2\%$ deviation is negative, placing $\ell_A$ in the descent-dependent population.
The Gram eigenvalue deficit (Part~0 Supplement) closes descent-dependent deviations;
$\ell_A$ requires a larger correction than the sub-2\% quantities, suggesting it is
more sensitive to the multi-step cascade descent. The Bott $\Omega_m = 0.31150$ gives
$\ell_A = 293.0$ ($-2.9\%$), a modest improvement.
\end{remark}

\section{The Hubble Constant}

\subsection{The cascade's lapse function}

The cascade's lapse at dimension $d$ is ([1], Section~3):
\[
N(d) = \sqrt{\pi} \cdot R(d) = \sqrt{\pi} \cdot \frac{\Gamma((d+1)/2)}{\Gamma((d+2)/2)}.
\]
At $d = 4$: $N(4) = \sqrt{\pi} \cdot \Gamma(5/2)/\Gamma(3) = 3\pi/8 = 1.17810$.

\subsection{The lapse-corrected Friedmann equation}

The standard Friedmann equation for a flat universe is $H_0^2 = \rho_{\rm total}/(3M_{\rm Pl,red}^2)$.
The cascade's invariant $I$ is computed in cascade time; the lapse $N(4)$ converts to
proper time. The corrected equation is $H_0 = M_{\rm Pl,red}^{-1} N(4)^{-1}\sqrt{I/(3\Omega_\Lambda)}$.

\begin{theorem}[Hubble constant]\label{thm:H0}
\begin{equation}
H_0 = \frac{8\,M_{\rm Pl,red}}{3\pi}\sqrt{\frac{\pi\,I}{3(\pi-1)}} = 71.05~\text{km/s/Mpc}.\label{eq:H0}
\end{equation}
\end{theorem}

\begin{proof}
$I = 1.0996 \times 10^{-120}$, $\Omega_\Lambda = (\pi-1)/\pi = 0.68169$, $N(4) = 3\pi/8$,
$M_{\rm Pl,red} = 2.435 \times 10^{18}$~GeV. Steps:
\begin{align*}
\pi I/[3(\pi-1)] &= 5.377 \times 10^{-121},\\
\sqrt{5.377 \times 10^{-121}} &= 7.333 \times 10^{-61},\\
M_{\rm Pl,red} \times 7.333 \times 10^{-61} &= 1.785 \times 10^{-42}~\text{GeV},\\
H_0 = (1.785 \times 10^{-42}~\text{GeV})/N(4) &= 71.05~\text{km/s/Mpc}.\qedhere
\end{align*}
\end{proof}

\subsection{Derivation of the lapse correction from the descent metric}

\begin{theorem}[Lapse correction from the descent metric]\label{thm:lapse}
The cascade descent from $d = \infty$ to $d = 4$ defines a foliation with lapse function
$N(d) = V_{d+1}/V_d$. The physical Friedmann equation at $d = 4$ is
$H_0^2 = I\,M_{\rm Pl,red}^2/(3\,\Omega_\Lambda\,N(4)^2)$, giving $H_0 = 71.05$~km/s/Mpc.
\end{theorem}

\begin{proof}
The cascade has two distinct time-like structures.

\textbf{Intra-level time.} The slicing coordinate $t$ within each level gives the Lorentzian
cascade metric $ds^2 = -dt^2 + (1-t^2)\,d\sigma_{d-1}^2$ ([3], Section~14.4). At the equator
$t = 0$: $\dot{a} = 0$, $H = 0$. The 4D observer sees a momentarily static $S^3$. This
is not the physical Hubble expansion.

\textbf{Inter-level time (the descent).} Paper~II, Section~7.1 identifies time with the
cascade's slicing direction. Physical expansion corresponds to sequential compactification
from $d = 217$ to $d = 4$.

In the ADM decomposition of the descent:
\[
ds^2_{\rm descent} = -N(d)^2\,dd^2 + a^2(d)\,d\sigma_3^2,
\]
where $N(d) = \sqrt{\pi}\cdot R(d)$ is the lapse, with $d\tau = N(d)\,dd$. The vacuum
energy density is $\rho_\Lambda = I \cdot M_{\rm Pl,red}^4$. The Friedmann equation in cascade
time gives $H_{\rm cascade}^2 = I\,M_{\rm Pl,red}^2/(3\,\Omega_\Lambda)$. Converting to proper time
via $H_{\rm phys} = H_{\rm cascade}/N(4)$:
\[
H_0 = \frac{8\,M_{\rm Pl,red}}{3\pi}\sqrt{\frac{\pi\,I}{3(\pi-1)}} = 71.05~\text{km/s/Mpc}.\qquad\square
\]
\end{proof}

\begin{remark}[Why the Gauss--Codazzi route does not produce the lapse]
The Gauss--Codazzi equations for the embedding chain are identically satisfied:
$K_{\mu\nu} = 0$ (Theorem~14.6 of [3]) and the Hamiltonian constraint telescopes to
$0 = 0$ (Corollary~14.7 of [3]). These confirm geometric self-consistency but do not
source the Friedmann equation. The matter content enters via the Bott partition (topological);
the lapse enters through the descent metric, not the intra-level Gauss--Codazzi equations.
\end{remark}

\begin{remark}[Self-consistency with the cosmological table]
The identification $\rho_\Lambda/M_{\rm Pl,red}^4 = I/N(4)^2 = 7.9 \times 10^{-121}$ is exactly
the Friedmann equation of Theorem~\ref{thm:lapse} solved for $\rho_\Lambda$. The observed
value $7.2 \times 10^{-121}$ deviates by $11\%$, consistent with the descent-dependent
systematic from subleading cascade corrections.
\end{remark}

\section{The Age of the Universe}

\begin{theorem}[Hubble time]
$1/H_0 = 13.76$~Gyr.
\end{theorem}

Numerically: $977.8~{\rm Gyr\cdot km/s/Mpc}/71.05 = 13.76$~Gyr.
Observed: $t_0 = 13.80 \pm 0.02$~Gyr (Planck 2018 [6]). Deviation: $-0.29\%$.

\begin{center}
\begin{tabular}{lcc}
\hline
Measurement & $H_0$ (km/s/Mpc) & Deviation from cascade \\
\hline
Planck 2018 (CMB) & $67.4 \pm 0.5$ & $-5.1\%$ \\
SH0ES (distance ladder) & $73.0 \pm 1.0$ & $+2.7\%$ \\
Cascade & 71.05 & --- \\
\hline
\end{tabular}
\end{center}

\section{The DESI BAO Observations}

\subsection{The sound horizon from cascade parameters}

The BAO standard ruler is the sound horizon at the drag epoch:
\[
r_d \approx 147.60\left(\frac{\omega_m}{0.1432}\right)^{-0.255}
\left(\frac{\omega_b}{0.02237}\right)^{-0.127}~\text{Mpc},
\]
where $\omega_x = \Omega_x h^2$ (Eisenstein \& Hu 1998 [12]). With cascade parameters
$H_0 = 71.05$, $h = 0.7105$:
\[
\omega_b^{\rm cas} = \tfrac{1}{2\pi^2} \times 0.7105^2 = 0.02557,\qquad
\omega_m^{\rm cas} = \tfrac{1}{\pi} \times 0.7105^2 = 0.16073.
\]

\begin{theorem}[Cascade sound horizon]
$r_d^{\rm cas} \approx 141$~Mpc, approximately $4.5\%$ smaller than the Planck-inferred
value $r_d^{\rm Planck} = 147.6$~Mpc.
\end{theorem}

\begin{proof}
$r_d = 147.60 \times (0.16073/0.1432)^{-0.255} \times (0.02557/0.02237)^{-0.127}
\approx 147.60 \times 0.971 \times 0.983 \approx 141$~Mpc.
\end{proof}

The difference has a clean physical origin: $\Omega_b = 1/(2\pi^2) = 0.0507$ exceeds
Planck's 0.0493, increasing baryon drag and shortening the sound horizon; simultaneously,
$H_0 = 71.05$ raises $h$ and further compresses $r_d$ through $\omega_b = \Omega_b h^2$.

\subsection{BAO observables and comparison to DESI DR2}

The table below is auto-generated from the cascade parameters via
\texttt{tools/generate\_bao\_table.py} and reproduced on every build. The
``Cascade'' column uses the leading-order $\Omega_m = 1/\pi$
($r_d^{\rm cas} \approx 141$~Mpc); the ``Bott'' column uses the
subleading-corrected $\Omega_m^{\rm Bott} = 0.31150$
($r_d^{\rm Bott} \approx 142$~Mpc). Both use $H_0 = 71.05$~km/s/Mpc.

% Auto-generated table: run tools/generate_bao_table.py before LaTeX build
%% Auto-generated by tools/generate_bao_table.py
% Cascade: H0=71.05, Omega_m=1/pi=0.31831, r_d=140.9 Mpc
% Bott:    H0=71.05, Omega_m=0.31150, r_d=141.7 Mpc
% Planck:  H0=67.4, Omega_m=0.315, r_d=147.6 Mpc
\begin{center}
\small
\begin{tabular}{llccccrr}
\hline
$z_{\rm eff}$ & Type & DESI DR2 & $\sigma$ & Cascade & Bott & Pull$_{\rm cas}$ & Pull$_{\rm Planck}$ \\
\hline
0.295 & $D_V/r_d$ & 7.930 & 0.150 & 7.967 & 7.938 & $+0.24\sigma$ & $+0.63\sigma$ \\
0.510 & $D_M/r_d$ & 13.650 & 0.200 & 13.346 & 13.305 & $-1.52\sigma$ & $-1.01\sigma$ \\
0.510$^\dagger$ & $D_H/r_d$ & 20.890 & 0.490 & 22.456 & 22.439 & $+3.20\sigma$ & $+3.60\sigma$ \\
0.706 & $D_M/r_d$ & 16.970 & 0.240 & 17.493 & 17.451 & $+2.18\sigma$ & $+2.76\sigma$ \\
0.706 & $D_H/r_d$ & 20.270 & 0.470 & 19.906 & 19.916 & $-0.77\sigma$ & $-0.38\sigma$ \\
0.934 & $D_M/r_d$ & 21.790 & 0.240 & 21.730 & 21.692 & $-0.25\sigma$ & $+0.50\sigma$ \\
0.934 & $D_H/r_d$ & 17.730 & 0.300 & 17.330 & 17.359 & $-1.33\sigma$ & $-0.76\sigma$ \\
1.321 & $D_M/r_d$ & 27.680 & 0.500 & 27.732 & 27.709 & $+0.10\sigma$ & $+0.58\sigma$ \\
1.321 & $D_H/r_d$ & 13.850 & 0.340 & 13.865 & 13.907 & $+0.05\sigma$ & $+0.47\sigma$ \\
1.484 & $D_M/r_d$ & 30.420 & 0.660 & 29.895 & 29.878 & $-0.80\sigma$ & $-0.40\sigma$ \\
1.484 & $D_H/r_d$ & 13.240 & 0.450 & 12.695 & 12.738 & $-1.21\sigma$ & $-0.91\sigma$ \\
2.330 & $D_M/r_d$ & 39.200 & 0.560 & 38.671 & 38.690 & $-0.94\sigma$ & $-0.30\sigma$ \\
2.330 & $D_H/r_d$ & 8.540 & 0.140 & 8.488 & 8.525 & $-0.37\sigma$ & $+0.31\sigma$ \\
\hline
\multicolumn{4}{l}{$\chi^2$ total} & 22.91 & 22.04 & & 24.65 \\
\multicolumn{4}{l}{$\chi^2/n$ ($n = 13$)} & 1.762 & 1.695 & & 1.896 \\
\hline
\end{tabular}
\end{center}
\noindent$^\dagger$: The $z = 0.510$, $D_H/r_d$ bin is a shared outlier; see Remark~\ref{rem:outlier}.


\subsection{The DESI dynamical dark energy signal as a ruler mismatch}

\begin{theorem}[Alternative account of the DESI signal]
The cascade cosmology ($w = -1$, $r_d \approx 141$~Mpc) fits DESI DR2 BAO with
$\chi^2/n = 1.70$ (Bott $\Omega_m$) or $1.76$ (leading $\Omega_m = 1/\pi$),
compared to Planck $\Lambda$CDM at $\chi^2/n = 1.90$. The apparent preference for dynamical
dark energy arises when Planck's $r_d = 147.6$~Mpc is imposed as an external prior;
the cascade's naturally smaller $r_d$ provides a $w = -1$ consistent alternative with
zero free parameters.
\end{theorem}

The logical structure: DESI directly observes the dimensionless ratios $D_M(z)/r_d$ and
$D_H(z)/r_d$. Converting to a preference for $w \neq -1$ requires an independent
determination of $r_d$---typically supplied by Planck CMB. The cascade predicts different
$(\Omega_b, H_0)$ values, hence a different $r_d$, and with that ruler the DESI distance ratios
are consistent with $w = -1$ throughout.

\begin{remark}[The $r_d$-independent check]
The ratio $D_M(z)/D_H(z)$ is independent of $r_d$ and tests the expansion geometry
directly. The cascade and Planck $\Lambda$CDM produce nearly identical predictions for this
ratio, indicating both models face the same geometric tension---driven by the $z = 0.510$
outlier---independently of ruler calibration.
\end{remark}

\subsection{Predicted apparent $w$ from the wrong ruler}

The cascade's true BAO distances (with $w = -1$, $r_d \approx 141$~Mpc) fitted by Planck's
standard pipeline (which assumes $r_d = 147.6$~Mpc) yield:
\begin{equation}
w_{\rm apparent} \approx -0.80.\label{eq:wapp}
\end{equation}

\begin{theorem}[Decomposition of the DESI signal]
The DESI DR2 dynamical dark energy signal decomposes into two independent components:
\begin{enumerate}
\item \textbf{The BAO-driven $w_0$ shift.} DESI BAO+CMB alone reports
$w_0 \approx -0.76 \pm 0.11$ [7]. The cascade ruler mismatch predicts
$w_{\rm apparent} \approx -0.80$. These agree within $1\sigma$. The $w_0$ component
is fully explained by the ruler error, with zero free parameters.
\item \textbf{The SNe-driven $w_a$ evolution.} $w_a \approx -0.6$ to $-1.1$ appears only
when Type~Ia supernova datasets are added. It is absent with Pantheon+ [10], and
its significance ranges from $2.5\sigma$ to $3.9\sigma$ depending on SNe calibration.
The cascade makes no prediction about SNe calibration; this component is neither
explained nor contradicted by the ruler mismatch.
\end{enumerate}
\end{theorem}

\begin{remark}[The $z = 0.510$ outlier]\label{rem:outlier}
The bin $z_{\rm eff} = 0.510$, $D_H/r_d$, shows a $3.20\sigma$ pull in the cascade and
$3.60\sigma$ in Planck $\Lambda$CDM. This single bin contributes $\approx 10.2$ of the
cascade's total $\chi^2 = 22.9$, and $\approx 13.0$ of Planck's total 24.7. Without it,
both fits would have $\chi^2/n \approx 1.0$. The tension persists in the $r_d$-free ratio
$D_M/D_H$---independent of ruler calibration---and represents a genuine anomaly in the
LRG1 expansion rate at $z \sim 0.5$. Neither model explains it.
\end{remark}

\subsection{Sensitivity of $r_d$ and $w_{\rm apparent}$ to cascade inputs}

\begin{center}
\begin{tabular}{lcccc}
\hline
$\Omega_m$ & $H_0$ & $r_d$ (Mpc) & $w_{\rm apparent}$ & $\Delta w$ \\
\hline
$0.299$ ($-4\%$) & 71.05 & 142.3 & $-0.838$ & $-0.051$ \\
$0.305$ ($-2\%$) & 71.05 & 141.7 & $-0.812$ & $-0.025$ \\
$1/\pi = 0.31831$ & 71.05 & $\approx141$ & $\approx-0.80$ & --- \\
$0.318$ ($+2\%$) & 71.05 & 140.3 & $-0.762$ & $+0.025$ \\
$0.324$ ($+4\%$) & 71.05 & 139.7 & $-0.738$ & $+0.049$ \\
\hline
\end{tabular}
\end{center}

A $\pm 4\%$ variation in $\Omega_m$ shifts $w_{\rm apparent}$ by $\pm 0.05$, remaining within
the DESI $1\sigma$ band. The prediction is robust to uncertainty in $\Omega_m$.

\subsection{Sharp falsification predictions}

Three nested tests:
\begin{enumerate}
\item \textbf{CMB-S4 measures $r_d$.} If $r_d \approx 141$~Mpc, the ruler mismatch is
confirmed and the DESI $w_0$ deviation dissolves. If $r_d > 144$~Mpc with
cascade-consistent $(H_0, \Omega_m)$, the cascade is falsified.
\item \textbf{DESI BAO+CMB $w_0$ stability.} The BAO-only $w_0$ should remain near
$-0.80$ while $w_a$ from BAO alone remains consistent with zero. If BAO alone drives
$w_a \neq 0$ at $> 3\sigma$, the ruler-mismatch explanation fails.
\item \textbf{SNe systematics.} If a recalibrated supernova dataset eliminates the $w_a$
signal, the remaining BAO-only $w_0 \approx -0.80$ is precisely what the cascade predicts.
\end{enumerate}

\section{Summary: The Complete Cosmological Table}

\begin{center}
\small
\begin{tabular}{llccc}
\hline
Parameter & Cascade formula & Predicted & Observed & Dev \\
\hline
$\rho_\Lambda/M_{\rm Pl}^4$ & $I/N(4)^2$ & $7.9\times10^{-121}$ & $7.2\times10^{-121}$ & $11\%$ \\
$\Omega_k$ & $\approx0$ ($S^3$ topology) & $\approx0$ & $0.001\pm0.002$ & --- \\
$\Omega_m$ & $1/\pi$ (lapse identity) & 0.31831 & $0.315\pm0.007$ & $+1.1\%$ \\
$\Omega_m^{\rm Bott}$ & Bott partition & 0.31150 & $0.315\pm0.007$ & $-1.1\%$ \\
$\Omega_\Lambda$ & $(\pi-1)/\pi$ & 0.68169 & $0.685\pm0.007$ & $-0.5\%$ \\
$\Omega_b$ & $1/(2\pi^2)$ & 0.05066 & $0.0493\pm0.0003$ & $+2.8\%$ \\
$\Omega_{\rm DM}$ & $(2\pi-1)/(2\pi^2)$ & 0.26765 & $0.264\pm0.007$ & $+1.4\%$ \\
$\Omega_r$ & $1/(4\pi^7)$ & $8.277\times10^{-5}$ & $8.27\times10^{-5}$ & $+0.1\%$ \\
$z_{\rm eq}$ & $4\pi^6$ & 3846 & $3402\pm26^*$ & --- \\
$T_{\rm CMB}$ & derived & 2.730~K & 2.7255~K & $+0.2\%$ \\
$w$ & $-1$ (fixed $\Lambda$, GB vanishes) & $-1$ & $-1.03\pm0.03$ & $1\sigma$ \\
$H_0$ & Eq.~\eqref{eq:H0} & 71.05 & 67.4--73.0 & between \\
$1/H_0$ & from $H_0$ & 13.76~Gyr & $13.80\pm0.02$ & $-0.3\%$ \\
$r_d$ & from $(\Omega_b,\Omega_m,H_0)$ & $\approx141$~Mpc & 141--148 & --- \\
$w_{\rm apparent}$ & ruler mismatch & $\approx-0.80$ & $-0.76\pm0.11$ & $<1\sigma$ \\
\hline
\end{tabular}
\end{center}
\noindent$^*$Planck's $z_{\rm eq}$ is model-dependent; not directly comparable.

\medskip All quantities are functions of $\pi$, the cascade invariant $I$, and the reduced
Planck mass $M_{\rm Pl,red}$ (the single dimensionful input). Zero fitted parameters.

\section{What This Paper Does Not Do}

\begin{itemize}
\item \textbf{Perturbation parameters.} The scalar spectral index $n_s$, primordial amplitude
$A_s$, and reionisation optical depth $\tau$ are not derived.
\item \textbf{The CMB power spectrum.} Requires $n_s$, $A_s$, $\tau$, and Boltzmann evolution
through recombination.
\item \textbf{The $z = 0.510$ anomaly.} Shared with Planck $\Lambda$CDM and not explained.
\end{itemize}

\section{Predictions and Falsifiability}

\begin{enumerate}
\item $w = -1$ \textbf{and} $r_d \approx 141$~Mpc (linked prediction). Three nested tests
follow (Section~8.5).
\item $\Omega_m = 1/\pi = 0.31831$ (leading order) / $\Omega_m^{\rm Bott} = 0.31150$ (subleading).
Falsified if future CMB measurements place $\Omega_m$ at $> 3\sigma$ from \emph{both}
values simultaneously.
\item $\Omega_b = 1/(2\pi^2) = 0.0507$. Currently $+2.8\%$ above Planck central value.
CMB-S4 will tighten to sub-percent; reanalysis with cascade parameters may shift the
inferred $\Omega_b$.
\item $H_0 = 70$--$72$~km/s/Mpc. Independent $H_0$ measurements (gravitational wave
standard sirens, time-delay cosmography) should converge on this range, not on either
extreme of the Hubble tension.
\item \textbf{No dark matter particles.} Confirmed detection of a dark matter particle
interacting via gauge forces would falsify the cascade's geometric dark matter
interpretation.
\end{enumerate}

\section{Open Questions}

\begin{enumerate}
\item \textbf{Theorem status of $\Omega_m$.} The Bott partition (Theorem~\ref{thm:omega_bott})
identifies non-trivial-phase layers with matter via the physical identification
hypothesis. Whether this identification step follows uniquely from the hypothesis,
or requires an additional argument, is the remaining open question for $\Omega_m$.
\item \textbf{Independent determination of $r_d$.} A direct CMB measurement without assuming
Planck's $(\Omega_b, H_0)$ would simultaneously test $\Omega_b = 1/(2\pi^2)$.
\item \textbf{Derive the cascade's expansion history $a(t)$.} The cascade's expansion history
should come from the descent profile $N(d)$, a function of the Gamma function.
\item \textbf{Perturbation spectrum from the cascade.} Whether the descent from $d = \infty$
generates density perturbations with a specific spectrum is the most important open
question in the cascade's cosmology programme.
\item \textbf{The $z = 0.510$ anomaly.} Not explained by the cascade or Planck $\Lambda$CDM. Open.
\end{enumerate}

\section{Confidence Assessment}

\textbf{Two-population systematic and the $\alpha(d^*)/\chi$ shift family.} At leading
order, descent-dependent quantities carry negative deviations and geometric quantities
carry positive deviations. Part~IVb identifies a structural family of cascade potential
shifts $\delta\Phi = \alpha(d^*)/\chi$ sourced at distinguished cascade layers from
Paper~I's four-dimension tower, all weighted by the same Euler characteristic
$\chi(S^{2n})=2$ that appears in the $2\sqrt{\pi}$ fermion obstruction. Two members of
the family close three SM precision observables at experimental precision:
\begin{itemize}
\item $\delta\Phi_{\mathrm{U}(1)}=\alpha(14)/\chi$ closes $\alpha_s(M_Z)=0.117917$ at
$+0.019\sigma$ and $m_\tau/m_\mu=16.81731$ at $+0.243\sigma$ together
(Part~IVb, Theorem~4.3);
\item $\delta\Phi_{\mathrm{phase}}=\alpha(19)/\chi$ closes the $\tau$ absolute mass at
$-0.31\sigma$, sourced at Paper~I's phase-transition threshold $d_1=19$
(Part~IVb, Theorem~4.4).
\end{itemize}
Both are closed forms in $\Gamma$ function values with zero fitted parameters. The
Part~0 Supplement's Gram-matrix first-order correction $\delta Q/Q_0 =
\sum_{\rm adj}(1-C_{d,d+1}^2)$ reduces the remaining descent deviations from $\sim\!2\%$
to sub-$1\%$: $v$ ($-1.0\%$), $\Omega_m^{\rm Bott}$ ($+1.0\%$). Geometric quantities
($m_H/m_W$, $\sin^2\theta_W$, $\theta_C$, $\Omega_b$, $\Omega_m = 1/\pi$, $\Omega_r$,
$T_{\rm CMB}$) are unchanged. $\Lambda$ at $+0.04\%$ belongs to neither population
(threshold quantity). Whether $v$, $m_W$, $\ell_A$, $\Omega_m$, and other remaining
descent-dependent quantities close at experimental precision via further members of
the $\alpha(d^*)/\chi$ family is an open question.

\textbf{Tier 1: Theorem-level.} $\Lambda = I$ (from [1]); $w = -1$ (fixed $\Lambda$,
Theorem~3.1); $\Omega_k \approx 0$ ($S^3$ topology); GB corrections vanish (two mechanisms,
Corollary~3.2); cascade fits DESI DR2 BAO with $w = -1$; ruler mismatch predicts apparent
$w \approx -0.80$, matching DESI BAO+CMB at $-0.76 \pm 0.11$ within $1\sigma$ (zero free
parameters); $H_0 = 71.05$~km/s/Mpc from the lapse-corrected Friedmann equation
(Theorem~\ref{thm:lapse}).

\textbf{Tier 2: Derived with clean argument.} $\Omega_m = 1/\pi$: lapse identity.
$\Omega_m^{\rm Bott} = 0.31150$: sphere-area fraction of non-trivial-phase layers, from the
Bott partition of [4] and Corollary~4.8a of [1] (Theorem~\ref{thm:omega_bott}).
$\Omega_b = 1/(2\pi^2)$: observer's boundary area (Theorem~\ref{thm:omega_b}).
$\Omega_r = 1/(4\pi^7)$: interior content, 7 geometric steps (Theorem~\ref{thm:omega_r}).
$H_0 = 71.05$~km/s/Mpc (Theorem~\ref{thm:H0}).

\textbf{What would change my mind.} (i) Robust $> 5\sigma$ measurement of $w \neq -1$ with
$r_d$ independently constrained at $> 144$~Mpc. (ii) Confirmed detection of a dark matter
particle. (iii) $\Omega_m$ measured at $> 3\sigma$ from \emph{both} $1/\pi = 0.31831$ and
the Bott value $0.31150$ simultaneously. (iv) CMB-S4 measurement of $r_d > 144$~Mpc with
$H_0$ consistent with the cascade. (v) Independent $H_0$ measurements converging outside
$68$--$74$~km/s/Mpc.

\begin{thebibliography}{12}

\bibitem{paper1} {[Author]}, \textit{Scale Variance from Orthogonality: How the Unit Ball
Generates $10^{120}$ Orders of Magnitude}, March 2026.

\bibitem{paper2} {[Author]}, \textit{Quantum Mechanics from the Cascade: Effective Theory of
a 4-Dimensional Observer in the Sphere-Area Geometry}, March 2026.

\bibitem{paper3} {[Author]}, \textit{General Relativity from the Cascade: Lovelock Uniqueness
and the Cosmological Constant of a 4-Dimensional Observer}, March 2026.

\bibitem{paper4a} {[Author]}, \textit{The Standard Model from the Cascade: Part IVa---Gauge
Group, Symmetry Breaking, and Three Generations from Bott Periodicity and Hairy Ball
Zeros}, March 2026.

\bibitem{paper4b} {[Author]}, \textit{The Standard Model from the Cascade: Part IVb---Masses,
Couplings, and Precision Predictions from the Geometric-Topological Factorization},
March 2026.

\bibitem{planck} Planck Collaboration, ``Planck 2018 results. VI. Cosmological parameters,''
\textit{Astronomy \& Astrophysics} \textbf{641}, A6 (2020).

\bibitem{desi2} DESI Collaboration, \textit{DESI DR2 Results II: Measurements of Baryon
Acoustic Oscillations and Cosmological Constraints}, arXiv:2503.14738 (2025).

\bibitem{desi1} DESI Collaboration, ``DESI 2024 VI: Constraints on models of dark energy,''
arXiv:2404.03002 (2024).

\bibitem{lee} S.~Lee, ``Comparing $\Lambda$CDM, $w$CDM, and $w_0w_a$CDM models with DESI DR2
BAO,'' arXiv:2507.01380 (2025).

\bibitem{riess} A.~G.~Riess et al., ``A comprehensive measurement of the local value of the
Hubble constant,'' \textit{Astrophysical Journal Letters} \textbf{934}, L7 (2022).

\bibitem{efstathiou} G.~Efstathiou, ``Evolving dark energy or supernovae systematics?''
\textit{Monthly Notices of the Royal Astronomical Society} \textbf{538}, 875 (2025).

\bibitem{EH98} D.~J.~Eisenstein and W.~Hu, ``Baryonic Features in the Matter Transfer
Function,'' \textit{Astrophysical Journal} \textbf{496}, 605 (1998).

\end{thebibliography}

\end{document}